%! suppress = TooLargeSection
%! suppress = SentenceEndWithCapital
%! suppress = TooLargeSection
% Preamble
\documentclass[11pt]{PyRollDocs}
\usepackage{textcomp}
\usepackage{csquotes}
\usepackage{wasysym}

\addbibresource{refs.bib}
% Document
\begin{document}

    \title{Karman Power and Labour PyRolL Plugin}
    \author{Christoph Renzing}
    \date{\today}

    \maketitle

    The PyRolL plugin pyroll-karman-power-and-labour calculates the roll-force and roll-torque as well as the mean neutral plane location for a roll-pass.
    Since von-Karman originally developed the strip or slap theory for flat rolling, the grooved pass is not considered directly.
    The solution is therefore derived for an equivalent flat pass, which makes the plugin's accuracy heavily depended on the equivalent flat pass plugins.

    The plugin provides three solvers, covering the full range from thick slabs to metal foils, and selects between them automatically for each roll pass:

    \begin{itemize}
        \item a rigid-roll solver based directly on von-Kármán's differential equation (\autoref{sec:model-approach}), valid whenever elastic effects are negligible;
        \item an elastic-plastic layer model (\autoref{sec:layer-model}), used for medium and thick passes, that adds elastic entry/exit zones, a mixed Coulomb/stiction friction law, through-thickness resolution into multiple layers, and full thermal coupling;
        \item a full elastic-contour foil-rolling model (\autoref{sec:foil-model}), used once elastic roll flattening becomes so severe that a circular roll-gap assumption is no longer valid.
    \end{itemize}

    \autoref{sec:dispatch} explains the criteria used to select between the three, and \autoref{sec:usage-instructions} lists the hooks the plugin defines.


    \section{Rigid-roll model (von-Kármán)}\label{sec:model-approach}

    The current model is based on the classic elementary theory of plasticity, which is a well known approach of modelling the stress distributions in flat rolling.
    The elementary theory of plasticity for flat rolling is also known as strip or slap theory, because the roll gap is divided in rolling direction in infinitesimal narrow strips, as shown in \autoref{fig:elementary-stripe}, on which the force balance is built.
    In thickness direction the strip's height is equal the local height of the roll gap.
    The deformation of each strip is always plane and parallelepidic.
    The strip's material properties are constant in thickness direction and may vary in rolling direction.
    Therefore, the vertical stress is constant in thickness direction and no shear deformation can be considered.

    \begin{figure}
        \centering
        \includegraphics[width=0.5\linewidth]{img/strip-hom}
        \caption{Stripe element of elementary theory of plasticity}
        \label{fig:elementary-stripe}
    \end{figure}

    The roll surface is described as a circular arc, since elastic roll deformation is neglected in this model (see \autoref{sec:layer-model} and \autoref{sec:foil-model} for the elastic-plastic and elastic-contour models, which lift this restriction).
    So the roll gap height can be described by \autoref{eq:roll-gap}.
    The rolls are here considered to be rigid.
    Also, the rolls are considered to be symmetric, meaning of equal radius $R_W$ and rotational frequency $n_W$.

    \begin{equation}
        h(x) = s + 2 * \left( R_W - \sqrt{R_W^2 - x^2} \right)
        \label{eq:roll-gap}
    \end{equation}

    The model needs additional approaches for the yield stress, and the friction between rolls and workpiece.
    The shear stress resulting from friction between the material and the roll surface is model by Coulomb's friction law.
    Since normally this model is used for cold rolling with lubricates, the friction coefficient $\mu$ has to be sufficiently high.
    The force balance at the stripe element leads to the well known Karmàn differential equation~\cite{Karman1925, Alexander1972} as shown in \autoref{eq:karman-elem}, which is an ordinary differential equation for the horizontal stress $\sigma_x$ in $x$ and can be solved easily by numerical methods.
    A suitable yield criterion gives the connection between $\sigma_x$ and $\sigma_y$.
    In this approach the yield criterion according to von-Mises was chosen~\cite{Mises1913, Mises1928}.

    \begin{gather}
        0 = \sigma_x, \frac{\mathrm{d}h}{\mathrm{d}x} + h\, \frac{\mathrm{d}\sigma_x}{\mathrm{d}x} + 2 \tau_R - 2 p_N \tan(\alpha)\\
        p_N = -\sigma_y - \tau_R \tan(\alpha)\\
        \tau = \mu p_N
        \label{eq:karman-elem}
    \end{gather}

    This rigid-roll solution never resolves elastic zones (the whole contact is assumed plastic from the first point of contact) and does not couple to temperature; it is implemented in \lstinline{karman_solver.py} as \lstinline{KarmanSolver}.


    \section{Elastic-plastic layer model for medium and thick passes}\label{sec:layer-model}

    For medium and thick passes, the rigid-roll assumption of \autoref{sec:model-approach} misses three effects that become significant as the relative pass thickness grows: an elastic entry/compression zone and elastic recovery zone at exit (rather than the whole contact being plastic immediately), a friction law that can saturate into sticking rather than remaining Coulomb throughout, and, once the strip is thick enough relative to the contact length, meaningfully different flow stress between the surface (which alone feels roll-contact friction) and the core through the pass. \lstinline{layer_solver.py} implements \lstinline{LayerRollingSolver}, a layer-count-parametrized model following Weiner~\cite{Weiner2021} (based on his \emph{RollingHomogeneousElasticPlastic} and \emph{RollingLayerModelElasticPlastic} implementations), which addresses all three:

    \begin{itemize}
        \item the strip is discretized into $N$ through-thickness layers (the single-layer case, $N=1$, is the "medium" regime; $N>1$ resolves the "thick slab" regime);
        \item each layer starts elastic and yields once its own local Mises/Tresca criterion is met, giving genuine elastic entry and exit zones per layer instead of a single sharp elastic/plastic split;
        \item friction at the roll interface, and between adjacent layers, follows a mixed Coulomb/sticking law after Bay and Wanheim~\cite{BayWanheim1976}, transitioning smoothly from Coulomb sliding to a sticking-limited shear stress as normal pressure grows;
        \item deformation heat, friction heat, conduction between neighbouring layers (and to the rolls), and convective transport are solved simultaneously with the stress equations, so each layer's temperature - and hence its flow stress - evolves individually through the pass.
    \end{itemize}

    Roll flattening is deliberately out of scope here, as it is for \lstinline{KarmanSolver} and \lstinline{FoilRollingSolver}'s own base radius: all three solvers simply read \lstinline{Roll.working_radius} once rather than deriving a flattened radius themselves. A separate plugin providing a Hitchcock~\cite{Hitchcock1935} (or other) flattened radius via that hook is therefore picked up transparently by every solver in this plugin, with no special-casing needed; the classical linear Hitchcock relation implemented in \lstinline{condition.py} (\lstinline{hitchcock_radius_ratio}) exists solely to \emph{classify} a pass for dispatch (\autoref{sec:dispatch}), and is never fed back into a calculation.

    Because each pass only sees the incoming profile's own scalar strain and temperature (\lstinline{pyroll-core} has no native concept of a through-thickness state carried from one pass to the next), all layers start uniform; the elastic entry zone therefore has no mechanism to break symmetry between layers (matching the reference implementation, which defines but never uses a redundant-shear correction after Ekelund~\cite{Ekelund1933}). Layers still diverge from the first plastification onward, since only the outermost layers feel roll-contact friction directly while inner layers feel the weaker inner-friction law, and the resulting temperature and flow-stress differences compound through the rest of the pass.


    \section{Full elastic-contour foil-rolling model}\label{sec:foil-model}

    Once elastic roll flattening becomes severe enough, even a scalar (Hitchcock) roll-flattening correction to \lstinline{Roll.working_radius} is no longer accurate, because the flattened roll-gap shape stops being well approximated by a single effective radius. \lstinline{foil_solver.py} implements \lstinline{FoilRollingSolver}, following Mauk and Overhagen~\cite{MaukOverhagen2013}, which instead computes the true, generally non-circular, deformed roll-gap shape directly: the normal-pressure distribution along the contact is superposed with elastic half-space (Boussinesq/Johnson~\cite{Johnson1985}) influence coefficients to obtain the roll surface's own elastic deflection, iterated to convergence with the strip's own elastic-plastic stress solution (elastic entry, plastic slip, an optional sticking zone, plastic slip, elastic exit).


    \section{Selecting between the three models}\label{sec:dispatch}

    \lstinline{roll_pass.py} wires the three solvers together behind a single hook, \lstinline{karman_solution}, so that \lstinline{roll_force}, \lstinline{Roll.roll_torque}, \lstinline{Roll.neutral_point} and the profile velocity hooks work unchanged regardless of which solver actually ran:

    \begin{enumerate}
        \item if the roll or profile elastic properties (\lstinline{elastic_modulus}, \lstinline{poissons_ratio}) are not set, the pass falls back to the rigid-roll model of \autoref{sec:model-approach} unconditionally, since none of the elastic effects can be evaluated;
        \item otherwise, the Hitchcock flattened-to-nominal radius ratio $r'/r$ is evaluated (Mauk and Overhagen~\cite{MaukOverhagen2013} give $r'/r \geq 2$ as the point where a circular roll-gap assumption is no longer trustworthy); if exceeded, the foil-rolling model of \autoref{sec:foil-model} is used;
        \item otherwise, the classical contact-length-to-mean-thickness ratio $L_d/H_m$ decides between the layer model's thick-slab (multiple layers) and medium (single layer) configurations, as described in \autoref{sec:layer-model}.
    \end{enumerate}

    All three limits are themselves hooks (\autoref{tab:hookspecs}) and can be overridden per pass, including forcing a particular model unconditionally.


    \section{Usage instructions}\label{sec:usage-instructions}
    The plugin can be loaded under the name \texttt{pyroll\_karman\_power\_and\_labour}.

    An implementation of the \lstinline{roll_force} hook on \lstinline{RollPass} is provided.
    Furthermore, an implementation of the \lstinline{roll_torque} and \lstinline{neutral_point} hooks on \lstinline{RollPass.Roll} is provided.

    Additionally, hooks on \lstinline{RollPass} are defined, which are used in the calculation, as listed in \autoref{tab:hookspecs}.
    The hooks \lstinline{back_tension}, \lstinline{front_tension} and \lstinline{coulomb_friction_coefficient} have to be set and adjusted individually; the elastic-plastic models additionally need \lstinline{elastic_modulus} and \lstinline{poissons_ratio} set on both \lstinline{RollPass.Roll} and \lstinline{RollPass.InProfile} to be selected at all (see \autoref{sec:dispatch}), and the layer model's thermal coupling needs \lstinline{specific_heat_capacity}, \lstinline{thermal_conductivity} and \lstinline{density} set on both.

    \begin{table}
        \centering
        \caption{Hooks specified by this plugin.}
        \label{tab:hookspecs}
        \begin{tabular}{@{}p{0.37\linewidth}p{0.55\linewidth}@{}}
            \toprule
            Hook name & Meaning \\
            \midrule
            \texttt{coulomb\_friction\_coefficient} & Coulomb's friction coefficient $\mu$ \\
            \texttt{back\_tension} & Back tension of the roll pass $\sigma_{x,0}$ \\
            \texttt{front\_tension} & Front tension of the roll pass $\sigma_{x,1}$ \\
            \texttt{karman\_solution} & Solver object (of one of the three models) with solution values as attributes \\
            \texttt{foil\_rolling\_condition} & Whether the foil-rolling model (\autoref{sec:foil-model}) is used \\
            \texttt{foil\_rolling\_hitchcock\_limit} & Hitchcock ratio $r'/r$ above which the foil model is used (default 2.0) \\
            \texttt{thick\_slab\_condition} & Whether the multi-layer, thick-slab configuration is used \\
            \texttt{layer\_model\_ld\_hm\_limit} & $L_d/H_m$ ratio below which the multi-layer configuration is used (default 1.0) \\
            \texttt{layer\_model\_layer\_count} & Number of through-thickness layers on the thick-slab path (default 5) \\
            \texttt{friction\_stiction\_coefficient} & Sticking friction coefficient feeding the Bay/Wanheim conversion (default 0.8) \\
            \texttt{roll\_heat\_transfer\_coefficient} & Heat transfer coefficient at the roll/material interface (default \SI{6000}{W/(m^2 K)}) \\
            \bottomrule
        \end{tabular}
    \end{table}

    \printbibliography


\end{document}
